The file manager has a web interface and a command-line tool on one core library. It encrypts all files with AES-256-GCM, protects keys with Argon2id, and uses HTTPS, CSRF tokens, a strict CSP, bound SQL parameters and owner checks on every file operation. I tested these controls with 53 automatic tests and with attacks on a running server. All tests passed, and all attacks failed.

The application has limits. The server sees the plain file during use, because the encryption runs on the server. Memory wiping and secure deletion are best-effort measures: Python and the library can keep copies, and SSDs, snapshots and backups can keep old data. The key cache, the pre-login tokens and the rate limit are in the memory of one process, so the server must run with one worker. The rate limit works per client address and has no lockout for one account. There is no multi-factor authentication and no password recovery. The application checks only the extension of an upload, and it does not scan files for malware. Behind a reverse proxy, the trusted proxy addresses and the rate limits must be configured, and request bodies, cookies and query strings must not be logged.

I learned three main lessons. First, a shared core library keeps the ownership checks the same in the web layer and in the CLI. Second, encrypted data must be checked before it is used, and error paths need as much care as normal paths: many bugs were in cleanup code after an error. Third, a check before opening a file is not enough, and the code must check the file that it really opened. The main challenges were to keep plain data in memory small for files up to 50 MiB, to update the libraries without changing the web behavior, and to test deletion and cleanup.
